\section{Lerch deformation: loss of real zeros and forced bifurcations}
\label{cont26:sec:lerch-bifurcations}

The manuscript proves eventual saturation of the positive-zero bound as
the parameter-derivative order grows, uniformly in the Lerch deformation
\cite{ProveIt}. At a fixed small derivative order the geometry is different.
We prove an exact count for sufficiently small positive deformation
parameters. Multiple elementary roots generally leave the real axis.
At indices three and four this forces an interior change in the number
of positive zeros before the Stieltjes endpoint is reached.

For integers $n\geq0$, $k\geq1$, let
\begin{equation}\label{cont26:eq:bif-family}
 f_{n,k}(a)=\frac{d^k}{da^k}\frac{\log^n a}{a},\qquad
 D_{n,k}^{\rho}(a)=\sum_{m=0}^{\infty}\rho^m f_{n,k}(a+m)
 \quad (a>0,\ 0\leq\rho\leq1).
\end{equation}
Here $0^0=1$ in the $m=0$ summand. The differentiated series is absolutely
convergent even at $\rho=1$, and
$D_{n,k}^{1}(a)=\gamma_n^{(k)}(a)$ for the Laurent convention
\[
 \zeta(s,a)=\frac1{s-1}
       +\sum_{n\geq0}\frac{(-1)^n}{n!}\gamma_n(a)(s-1)^n.
\]
For $\rho<1$, it equals $\partial_a^k\ell_n(\rho,a)$ with
$\ell_n(\rho,a)=\sum_{m\geq0}\rho^m\log^n(a+m)/(a+m)$.
In particular, these are order derivatives of the Lerch transcendent;
at $a=1$ their undeformed-in-$a$ antecedents specialize to order
derivatives of $\operatorname{Li}_s(\rho)$.

Define the rational polynomial
\begin{equation}\label{cont26:eq:bif-polynomial}
 P_{n,k}(x)=n![t^n]\left(e^{xt}\prod_{j=1}^k(1+t/j)\right).
\end{equation}
Coefficient extraction after differentiating $a^{-1-t}$ gives
\begin{equation}\label{cont26:eq:bif-elementary}
 (-1)^{n+k} f_{n,k}(a)=k!a^{-k-1}P_{n,k}(-\log a).
\end{equation}
The sign normalization
$F_{n,k}^{\rho}=(-1)^{n+k}D_{n,k}^{\rho}$ does not affect zeros.

\subsection{The elementary roots and uniform exclusion of escaping zeros}

\begin{lemma}\label{cont26:lem:bif-elementary-roots}
If $n>k\geq1$ and $d=n-k$, then $P_{n,k}$ has a root of multiplicity
$d$ at zero and precisely $k$ other roots, all simple and negative.
Its first nonzero coefficient is
\[
 [x^d]P_{n,k}(x)=\frac{n!}{d!\,k!}.
\]
\end{lemma}
\begin{proof}
The polynomial is $\prod_{j=1}^k(1+\partial_x/j)x^n$. If a polynomial
$p$ has real roots $r_i$ of multiplicities $m_i$, the new noninherited
roots of $p+c p'$ solve
\[
 \frac{p'(x)}{p(x)}=\sum_i\frac{m_i}{x-r_i}=-1/c.
\]
For $c>0$ the left side decreases strictly between its poles. There is
one new simple root in every intervening interval and one to the left
of all the old roots, and there is none to their right. Every inherited
multiple root loses one unit of multiplicity. Induction from $x^n$
therefore gives the asserted configuration. The coefficient follows
from the top coefficient $1/k!$ in $\prod_{j=1}^k(1+t/j)$.
\end{proof}

Let $A_{n,k}>1$ be the largest zero of $f_{n,k}$, equivalently the
exponential of the negative of the smallest root of $P_{n,k}$.

\begin{lemma}[No escape through either endpoint]\label{cont26:lem:bif-no-escape}
For fixed $n>k\geq1$ there is $\varepsilon_{n,k}>0$ such that every
positive zero of $D_{n,k}^{\rho}$, for every $\rho\in[0,1]$, lies in
$[\varepsilon_{n,k},A_{n,k}]$. For $\rho>0$ the upper inequality is strict.
Moreover $D_{n,k}^{\rho}\to f_{n,k}$ uniformly on every compact
subinterval of $(0,\infty)$ as $\rho\downarrow0$, including every fixed
number of $a$-derivatives.
\end{lemma}
\begin{proof}
The $m=0$ term in the normalized series satisfies
\[
 F_{n,k}^{0}(a)\sim k!a^{-k-1}\log^n(1/a)\longrightarrow+\infty
 \qquad(a\downarrow0).
\]
For $0<a\leq1$, the sum of the absolute values of all $m\geq1$ terms is
bounded independently of $\rho\in[0,1]$: each is bounded by a constant
times $(1+\log^n(m+1))/m^{k+1}$, whose series converges.
Thus the $m=0$ term dominates uniformly near zero.

For $a>A_{n,k}$, the argument $-\log(a+m)$ lies strictly to the left
of every root of $P_{n,k}$ for all $m\geq0$. Consequently all summands
have the same nonzero sign. At $a=A_{n,k}$ the $m=0$ summand vanishes
and the others have that common strict sign when $\rho>0$. This excludes
the upper endpoint and infinity. On a compact positive $a$-interval,
the same summable majorants, with $k$ increased for the derivatives,
give the final assertion by dominated convergence.
\end{proof}

\subsection{Complete small-parameter count and Puiseux branches}

\begin{theorem}[Unfolding the multiple elementary root]
\label{cont26:thm:bif-small-rho}
Fix $n>k\geq1$, put $d=n-k$, and define
\begin{equation}\label{cont26:eq:bif-constants}
 A=\frac{n!}{d!},\qquad
 B=\frac{k!}{2^{k+1}}P_{n,k}(-\log2).
\end{equation}
Assume $B\ne0$. For all sufficiently small $\rho>0$, every positive
zero of $D_{n,k}^{\rho}$ is simple and their exact number is
\begin{equation}\label{cont26:eq:bif-count}
 k+\begin{cases}
 1,&d\text{ odd},\\
 2,&d\text{ even and }B<0,\\
 0,&d\text{ even and }B>0.
 \end{cases}
\end{equation}
The $k$ zeros continuing the elementary zeros different from one are
real analytic functions of $\rho$. All $d$ complex zeros near $a=1$
are convergent Puiseux branches. Writing $t=\rho^{1/d}$ and letting
$y_1,\ldots,y_d$ be the distinct roots of
\begin{equation}\label{cont26:eq:bif-leading-equation}
 A y^d+B=0,
\end{equation}
these branches have the form
\begin{equation}\label{cont26:eq:bif-puiseux}
 a_j(\rho)=1-y_j\rho^{1/d}+O(\rho^{2/d}).
\end{equation}
Their series converge for sufficiently small $t$. In particular,
if $d\geq3$, the number in \eqref{cont26:eq:bif-count} is strictly less than $n$.
\end{theorem}
\begin{proof}
For $|\rho|<1$ the series is normally convergent on compact subsets
of $\Re a>0$ and hence jointly holomorphic in $(a,\rho)$. The principal
logarithm is unambiguous on this half-plane. The implicit-function
theorem continues each of the $k$ nonunit elementary zeros uniquely
and real analytically.

Near the remaining zero use the analytic coordinate $a=e^{-x}$.
Equations~\eqref{cont26:eq:bif-elementary} and \eqref{cont26:eq:bif-constants} give
\begin{equation}\label{cont26:eq:bif-local}
 F_{n,k}^{\rho}(e^{-x})
       =A x^d+B\rho+O(x^{d+1}+\rho x+\rho^2).
\end{equation}
Substitute $\rho=t^d$ and $x=ty$, divide by $t^d$, and extend
analytically to $t=0$. The resulting function equals $Ay^d+B$ at zero.
Its $d$ roots are simple because $AB\ne0$, so the implicit-function
theorem supplies $d$ analytic branches $y_j(t)=y_j+O(t)$. Rouch\'e's
theorem on a fixed small circle around $x=0$ shows that these exhaust
the zeros converging to zero. A branch is real for small positive $t$
exactly when its leading root $y_j$ is real: real leading roots continue
by conjugation and uniqueness; nonreal leading roots stay nonreal.
The polynomial $Ay^d+B$ has precisely the real-root count displayed
in \eqref{cont26:eq:bif-count}, since $A>0$.

Finally, Lemma~\ref{cont26:lem:bif-no-escape} confines every positive zero to
a fixed compact interval. Delete disjoint small neighborhoods of the
elementary zeros from that interval. The limit $f_{n,k}$ is bounded
away from zero on the remaining compact set; uniform convergence
excludes every additional zero there. Thus the local branch count is
the global positive-zero count. Simplicity follows from the nonzero
implicit derivatives, at both the ordinary and rescaled branches.
\end{proof}

The nonvanishing condition is displayed explicitly in the general
statement. For the following first-derivative specialization, its sign
and nonvanishing have an elementary proof for every index.

\begin{corollary}[All indices at first derivative order]
\label{cont26:cor:bif-kone}
For every fixed integer $n\geq2$, and all sufficiently small $\rho>0$,
$D_{n,1}^{\rho}$ has exactly one positive zero if $n$ is odd and
exactly two if $n$ is even. All are simple. One zero tends to $e^n$
and depends real analytically on $\rho$. When $n$ is even, the second
zero satisfies
\begin{equation}\label{cont26:eq:bif-even-root}
 a_{\mathrm{small}}(\rho)=1-\kappa_n\rho^{1/(n-1)}
             +O(\rho^{2/(n-1)}),\qquad
 \kappa_n=\log2\left(\frac{n-\log2}{4n}\right)^{1/(n-1)}>0.
\end{equation}
The root tending to $e^n$ moves initially to the left:
\begin{align}\label{cont26:eq:bif-large-root}
 a_{\mathrm{large}}(\rho)&=e^n-C_n\rho+O(\rho^2),\\
 C_n&=\frac{e^{3n}\log^{n-1}(e^n+1)
                    (\log(e^n+1)-n)}{n^{n-1}(e^n+1)^2}>0.\notag
\end{align}
\end{corollary}
\begin{proof}
Here
\[
 P_{n,1}(x)=x^{n-1}(x+n),\qquad
 f_{n,1}(a)=\frac{\log^{n-1}a\,(n-\log a)}{a^2}.
\]
Since $0<\log2<1<n$, the constant $B$ has sign $(-1)^{n-1}$.
For odd $n$ the multiplicity $n-1$ is even and $B>0$, giving no
real branch near one. For even $n$ it is odd and $B<0$, giving one
positive leading root $y=\kappa_n$ in \eqref{cont26:eq:bif-leading-equation}.
This proves the count and \eqref{cont26:eq:bif-even-root}.
At $a=e^n$,
$f'_{n,1}(e^n)=-n^{n-1}e^{-3n}$. Implicit differentiation of
$f_{n,1}(a)+\rho f_{n,1}(a+1)+O(\rho^2)=0$ gives
$a'(0)=-f_{n,1}(e^n+1)/f'_{n,1}(e^n)=-C_n$.
\end{proof}

For example, the double elementary zero for $n=3$ gives a conjugate
nonreal pair with
$a=1\mp i\kappa_3\sqrt\rho+O(\rho)$, whereas for $n=4$ the triple zero
gives one real branch below one and a nonreal conjugate pair. The
manuscript's eventual simplicity theorem is fully consistent with
this result: its derivative order tends to infinity, while the present
theorem fixes the derivative order and lets the deformation tend to zero.

\subsection{The logarithmic singularity at the Stieltjes endpoint}

An elementary tail calculation supplies a complementary expansion as
$\rho\uparrow1$. It gives the precise regularity that is not visible
in continuity of the positive-kernel representation. In this subsection
the indices are fixed and the limit concerns the deformation parameter.

For $0\leq j<k$ define the absolutely convergent moments
\begin{equation}\label{cont26:eq:bif-weighted-moments}
 M_{n,k,j}(a)=\sum_{m=0}^{\infty}m^j f_{n,k}(a+m),
 \qquad M_{n,k,0}(a)=\gamma_n^{(k)}(a).
\end{equation}
The $m=0$ factor is one when $j=0$ and zero otherwise.

\begin{theorem}[Endpoint expansion with its first nonanalytic term]
\label{cont26:thm:bif-rho-one}
Fix integers $n\geq0$, $k\geq1$ and a compact interval
$I\subset(0,\infty)$. Put $\epsilon=-\log\rho$ and
$L=\log(1/\epsilon)$. As $\epsilon\downarrow0$, uniformly for $a\in I$,
\begin{align}\label{cont26:eq:bif-rho-one}
 D_{n,k}^{e^{-\epsilon}}(a)
  ={}&\sum_{j=0}^{k-1}\frac{(-\epsilon)^j}{j!}M_{n,k,j}(a)
       +\frac{\epsilon^k L^{n+1}}{n+1}\\
     &\quad+O_{n,k,I}\bigl(\epsilon^k(1+L^n)\bigr).\notag
\end{align}
Thus the family is $C^{k-1}$ at the endpoint in the deformation
parameter, but its $k$th one-sided derivative is not finite there.
In particular it is not $C^k$. The coefficient of the leading
nonanalytic term is independent of $a$ and $k$.
\end{theorem}
\begin{proof}
For large $m$, uniformly on $I$, differentiation of
$\log^n(a+m)/(a+m)$ gives
\begin{equation}\label{cont26:eq:bif-large-m}
 f_{n,k}(a+m)=(-1)^k k!\frac{\log^n m}{m^{k+1}}+E_{n,k}(a,m).
\end{equation}
For $n\geq1$ the error is bounded by a constant times
\[
 \frac{1+\log^{n-1}m}{m^{k+1}}
       +\frac{1+\log^n m}{m^{k+2}};
\]
for $n=0$ it is $O(m^{-k-2})$. Indeed, the leading term arises when
all $k$ derivatives fall on $1/(a+m)$; every differentiation of a
logarithm lowers its degree, and replacement of $a+m$ by $m$ adds
one inverse power. This also proves convergence and local uniform
convergence of all the moments in \eqref{cont26:eq:bif-weighted-moments}.

Write
\[
 r_k(v)=e^{-v}-\sum_{j=0}^{k-1}\frac{(-v)^j}{j!}.
\]
Taylor's theorem and the finite polynomial bound give
\begin{align}\label{cont26:eq:bif-exponential-remainders}
 r_k(v)&=\frac{(-1)^k v^k}{k!}+O_k(v^{k+1})
                         &&(0\leq v\leq1),\\
 |r_k(v)|&\leq C_k v^{k-1}&&(v\geq1).\notag
\end{align}
Subtract the finite moment sum in \eqref{cont26:eq:bif-rho-one} from the
defining series. The remainder is exactly
$\sum_{m\geq1}f_{n,k}(a+m)r_k(\epsilon m)$.
Put $N=\lfloor1/\epsilon\rfloor$. In the range $m\leq N$ the error
from the first line of \eqref{cont26:eq:bif-exponential-remainders} is bounded by
\[
 C\epsilon^{k+1}\sum_{m\leq N}
                 m^{k+1}|f_{n,k}(a+m)|
   \leq C'\epsilon^k(1+L^n).
\]
Its retained part is
\[
 \frac{(-1)^k\epsilon^k}{k!}
       \sum_{m\leq N}m^k f_{n,k}(a+m).
\]
By \eqref{cont26:eq:bif-large-m}, ordinary integral comparison for logarithmic
harmonic sums gives
\[
 \sum_{m\leq N}m^k f_{n,k}(a+m)
   =(-1)^k k!\left\{\frac{L^{n+1}}{n+1}
                                +O(1+L^n)\right\}.
\]
For $n=0$ the error contributed by \eqref{cont26:eq:bif-large-m} is bounded;
for $n\geq1$ its lower logarithmic terms sum to $O(1+L^n)$.
The replacement of $\log N$ by $L$ is absorbed in the same error.
The two factors $(-1)^k$ cancel, producing the positive coefficient
in the theorem.

Finally the range $m>N$, using the second bound in
\eqref{cont26:eq:bif-exponential-remainders}, is at most
\[
 C\epsilon^{k-1}\sum_{m>N}\frac{1+\log^n m}{m^2}
       =O\bigl(\epsilon^k(1+L^n)\bigr).
\]
This proves the expansion with the claimed uniformity.

Differentiation with respect to $\epsilon$ through order $k-1$ is
justified at zero by the absolutely convergent weighted moment series,
uniformly on $I$. For clarity, the failure of the next derivative can
also be checked directly, without differentiating a big-$O$ term.
Put $G(\epsilon)=D_{n,k}^{e^{-\epsilon}}(a)$. Then
\[
 G^{(k-1)}(\epsilon)
   =(-1)^{k-1}\sum_{m\geq0}m^{k-1}e^{-\epsilon m}f_{n,k}(a+m).
\]
Apply the same split at $m=\lfloor1/\epsilon\rfloor$ to the difference
between this expression and its value at zero. Its summand has leading
coefficient $-k!\log^n m/m^2$, by \eqref{cont26:eq:bif-large-m}, and therefore
\[
 \frac{G^{(k-1)}(\epsilon)-G^{(k-1)}(0)}{\epsilon}
   =\frac{k!}{n+1}L^{n+1}+O(1+L^n)\longrightarrow+\infty.
\]
This proves the failure of a finite $k$th one-sided derivative. The
change of coordinates $\rho=e^{-\epsilon}$ is smooth with nonzero
derivative at zero; its lower-order chain-rule terms have finite limits.
Thus the same endpoint regularity holds in $\rho$.
\end{proof}

\begin{corollary}[Displacement of a simple Stieltjes zero]
\label{cont26:cor:bif-stieltjes-displacement}
Fix an integer $n\geq1$. Let $a_*>0$ be a simple zero of $\gamma_n'(a)$, so
$\gamma_n''(a_*)\ne0$. For $\rho<1$ sufficiently close to one there is a
unique nearby simple zero $a_*(\rho)$ of $D_{n,1}^{\rho}$. With
$\epsilon=-\log\rho$ and $L=\log(1/\epsilon)$,
\begin{equation}\label{cont26:eq:bif-root-displacement}
 a_*(\rho)=a_*
   -\frac{\epsilon L^{n+1}}{(n+1)\gamma_n''(a_*)}
   +O\bigl(\epsilon(1+L^n)\bigr).
\end{equation}
\end{corollary}
\begin{proof}
Local uniform convergence of the series and its $a$-derivative gives
a unique simple zero near $a_*$. The $k=1$ specialization of
Theorem~\ref{cont26:thm:bif-rho-one} is
\[
 D_{n,1}^{e^{-\epsilon}}(a)-\gamma_n'(a)
     =\frac{\epsilon L^{n+1}}{n+1}
           +O\bigl(\epsilon(1+L^n)\bigr),
\]
uniformly on a small interval around $a_*$. A derivative bounded away
from zero first gives $a_*(\rho)-a_*=O(\epsilon L^{n+1})$.
Taylor expansion of $\gamma_n'$ at $a_*$ then gives the displayed
formula. Its quadratic error is absorbed because
$\epsilon^2L^{2n+2}=o(\epsilon(1+L^n))$.
\end{proof}

The logarithmic enhancement means that a very small value of $1-\rho$
can still noticeably displace a fixed-index zero. This is a proved
endpoint scale. It does not determine where an interior multiple root
occurs, because the simple-root expansion ceases to apply when the
derivative in its denominator approaches zero.

\subsection{Fresh exact endpoint signs}

We recall, with proof, the global zero bound used to upgrade sign changes
to exact endpoint counts. It is the result entitled ``Global bound on
positive zeros'' in Chapter~10 of the source manuscript; the accompanying
results are entitled ``Strictly real-rooted reciprocal-gamma polynomials''
and ``Appell--Laplace representation'' \cite{ProveIt}.
Define reciprocal-gamma Appell polynomials $Q_n$ by
\[
 \frac{e^{xt}}{\Gamma(1+t)}=\sum_{n\geq0}Q_n(x)\frac{t^n}{n!}.
\]
\begin{lemma}[Global positive-zero bound]\label{cont26:lem:bif-global-count}
For every $n\geq0$, $k\geq1$, and $0\leq\rho\leq1$, the function
$D_{n,k}^{\rho}$ has at most $n$ positive zeros, counted with multiplicity.
\end{lemma}
\begin{proof}
The Weierstrass product is
\[
 \frac1{\Gamma(1+t)}=e^{\gamma t}
             \prod_{j\geq1}(1+t/j)e^{-t/j}.
\]
Taking $n![t^n]$ after multiplication by $e^{xt}$ turns finite products
into the operators $1+\partial_x/j$ and real translations applied to
$x^n$. These preserve real-rootedness by the operator proof in
Lemma~\ref{cont26:lem:bif-elementary-roots}. Local uniform convergence of the
product gives coefficient convergence of the resulting monic polynomials,
so the limiting $Q_n$ is real-rooted. For $n\geq1$, remove the first
$n-1$ factors before taking the limit. The remaining tail still gives
a real-rooted monic polynomial. Restoring the removed factors, including
their commuting real translations, applies $n-1$ root-preserving
operators, which reduce every inherited multiplicity until all roots
are simple. Thus $Q_n$ has exactly $n$ distinct real roots. For $n=0$
one has $Q_0=1$.

The Mellin integral gives the identity
\[
 \sum_{n\geq0}\frac{(-1)^{n+k}}{n!}D_{n,k}^{\rho}(a)t^n
   =\frac1{\Gamma(1+t)}\int_0^\infty
                   \frac{x^{k+t}e^{-ax}}{1-\rho e^{-x}}\,dx
\]
for $|t|<k$ in a neighborhood of zero. It follows either by
differentiating $(a+m)^{-1-t}$ $k$ times, or by canceling
$\Gamma(k+1+t)$ against its rising factorial. Coefficient comparison
proves
\begin{equation}\label{cont26:eq:bif-laplace}
 F_{n,k}^{\rho}(a)=\int_0^\infty
       \frac{x^ke^{-ax}}{1-\rho e^{-x}}Q_n(\log x)\,dx.
\end{equation}
At zero, $1/(1-\rho e^{-x})\leq1/(1-e^{-x})=O(x^{-1})$,
so $x^{k-1}|\log x|^n$ supplies an integrable majorant. At infinity
$e^{-ax}$ dominates all the logarithmic and polynomial factors.
These estimates justify coefficient extraction, differentiation in $a$,
and multiplication of the kernel by any fixed polynomial.

The signed kernel in \eqref{cont26:eq:bif-laplace} has exactly $n$ sign changes.
For completeness, let $G$ be a Laplace transform of such a kernel.
If there are no sign changes, $G$ has a fixed strict sign. Otherwise
choose one sign-change point $x_0$. Multiplication of the kernel by
$x_0-x$ removes that sign change and introduces none. Its transform is
$(\partial_a+x_0)G=e^{-x_0a}(e^{x_0a}G)'$ and, by induction, has at
most $n-1$ zeros counted with multiplicity and is not identically zero.
Therefore $G$ also cannot be identically zero. Any finite collection of
$M$ zeros of $G$, counted with multiplicity, forces at least $M-1$ zeros
of $(e^{x_0a}G)'$ by Rolle's theorem, including inherited multiplicities.
Hence $M\leq n$. The assertion holds for every finite collection and
therefore for all positive zeros, completing the induction and proof.
\end{proof}

The following are newly replayed rational enclosures for
\[
 V_n(a)=(-1)^{n+1}a^2\gamma_n'(a).
\]
Each displayed decimal endpoint is an exact terminating rational number.
\begin{center}
\begin{tabular}{cclc}
\hline
$n$&$a$&Interval containing $V_n(a)$&Sign\\
\hline
3&$4/5$&$[0.163799,\ 0.163919]$&$+$\\
3&$11/10$&$[-0.041066,\ -0.040839]$&$-$\\
3&$3/2$&$[0.055828,\ 0.056247]$&$+$\\
3&$21/10$&$[-0.245836,\ -0.245015]$&$-$\\
4&$4/5$&$[0.038668,\ 0.039282]$&$+$\\
4&$47/50$&$[-0.003812,\ -0.002965]$&$-$\\
4&$6/5$&$[0.021074,\ 0.022453]$&$+$\\
4&$17/10$&$[-0.037951,\ -0.035186]$&$-$\\
4&$11/5$&$[0.148603,\ 0.153231]$&$+$\\
\hline
\end{tabular}
\end{center}

For reproducibility, the included script uses only Python's standard
library and outward integer arithmetic on the grid $2^{-160}\mathbb Z$.
It sums the first $M=1000$ terms of \eqref{cont26:eq:bif-family} at $k=1,\rho=1$.
For a positive rational argument, power-of-two range reduction brings
the logarithm to $1\leq x\leq2$. With $u=(x-1)/(x+1)\in[0,1/3]$, it uses
\begin{equation}\label{cont26:eq:bif-log-certificate}
 0\leq\log x-2\sum_{j=0}^{J-1}\frac{u^{2j+1}}{2j+1}
 \leq\frac{2u^{2J+1}}{(2J+1)(1-u^2)},\qquad J=64.
\end{equation}
All powers, divisions and the geometric remainder are enclosed by
integer floor and ceiling. Inverting the argument handles logarithms
below one. Ordinary interval multiplication then evaluates $f_{n,1}$.

The tail needs only a monotone integral comparison. Put $b=a+M$ and
$h_n(x)=-f_{n,1}(x)$. For $n=3,4$ and $\log x\geq6$, this function is
positive and decreasing, since
\[
 h_n'(x)=-\frac{\log^{n-2}x}{x^3}
          \{2\log^2x-3n\log x+n(n-1)\}<0.
\]
Every argument $b$ above satisfies $\log b>6$. As
$\int_b^\infty h_n(x)\,dx=\log^n b/b$, one has the exact enclosure
\begin{equation}\label{cont26:eq:bif-tail-certificate}
 f_{n,1}(b)-\frac{\log^n b}{b}
 \leq\sum_{m=M}^{\infty}f_{n,1}(a+m)
 \leq-\frac{\log^n b}{b}.
\end{equation}
Combining these bounds proves every displayed sign; the JSON artifact
also contains the finer integer endpoints before decimal display rounding.
Thus these signs do not depend on quadrature or a special-function library.

\begin{corollary}[Forced interior multiple roots]
\label{cont26:cor:bif-interior}
For each $n\in\{3,4\}$ there are $\rho_*\in(0,1)$ and $a_*>0$ such that
\begin{equation}\label{cont26:eq:bif-discriminant}
 D_{n,1}^{\rho_*}(a_*)=0,
 \qquad \partial_aD_{n,1}^{\rho_*}(a_*)=0.
\end{equation}
At $\rho=1$ there are exactly $n$ positive zeros, all simple, and this
count persists for $\rho<1$ sufficiently close to one. For small positive
$\rho$ the counts are instead one and two, respectively.
\end{corollary}
\begin{proof}
The certified signs give $n$ disjoint intervals with sign changes.
The global bound with multiplicity shows that these are all the roots
and each is simple. Local uniform convergence as $\rho\uparrow1$ of the
series and its $a$-derivative preserves a unique simple root in each
small interval. The global bound excludes extras. The small-parameter
counts are Corollary~\ref{cont26:cor:bif-kone}.

If every root were simple throughout $0<\rho<1$, the implicit-function
theorem and compact exclusion of both endpoints would make the number
of positive roots locally constant on this connected interval, hence
constant. More explicitly, finitely many small root neighborhoods and
the zero-free compact complement yield a common parameter neighborhood
on which no root can be added or lost. The different endpoint counts
contradict constancy. Therefore \eqref{cont26:eq:bif-discriminant} holds at
an interior parameter.
\end{proof}

Ordinary numerical quadrature of \eqref{cont26:eq:bif-laplace}, followed by a
two-variable solve, suggests transitions near
\[
\begin{array}{c|cc}
n&a_*&\rho_*\\ \hline
3&1.0497530509&0.99999046684\\
4&1.2013268276&0.99999992971
\end{array}
\]
These coordinates are numerical observations. The theorem establishes
an interior multiple root; it does not establish uniqueness of the
transition, multiplicity exactly two, or a nonzero derivative in the
deformation parameter. Certifying those stronger statements is a
specific further problem, approachable by a two-variable interval
Newton calculation followed by an analytic exclusion argument.

